\chapter{生命孵化场}

地球在星图里放大。

它悬在众人之间，蓝白色光芒把导航室破损的墙面照亮。法厄同的警告声被压低，仿佛连濒死星球都暂时让位给这颗年轻行星。

那时的地球还没有我熟悉的轮廓。

海洋辽阔得像没有边界，陆地分布陌生，大气中翻滚着厚重云带。轨道上没有城市，没有卫星，也没有后世人类留下的任何痕迹。它只是安静地转着，像一枚尚未打开的种子。

但在各族眼里，它显然不是“只是”一颗行星。

星图自动展开多层标注。海洋区域浮出蓝色频谱，陆地边缘出现绿色适应曲线，大气层上方有一圈细密的意识谐振读数。那些读数我看不懂，洛安也看不懂，但房间里的其他族群都懂一些。

日辉族幸存者低声说：“抽象意识阈值尚未成形。”

鳞裔族女人接着说：“大型载体方向未锁定。”

渊民族老人看向海洋：“海已经能记忆，只是还不能说话。”

节肢族个体的节点闪烁：“集群潜力存在。个体隔离风险高。”

瑟恩没有评价。他只看着地球周围那些交错航线，脸色越来越难看。

查理忽然明白，这不是一张地图。

这是体检报告。

在他们的标注里，地球没有故乡或神赐之地的光环。它被拆成一个正在发育的生命孵化场。海洋、陆地、大气、重力、磁场、神经演化潜力、意识承载阈值，全部被测量、标注、比较。

我感到一阵说不出的愤怒。

不是因为他们邪恶。

恰恰是因为他们太认真。

他们像医生、工程师、园丁、设计师，围着一个还没有说话能力的孩子，争论这个孩子长大后应该成为谁。每一族都能说出理由，每一族都有数据，每一族都相信自己的方案能让未来少受苦。

可没有人问地球。

当然，地球也无法回答。

星图继续展开。

日辉族方案在地球陆地上投下金色网格。未来城市、祭坛、观测塔、王权中心和律令节点一一浮现。鳞裔族方案则沉入地壳，形成地下库、长眠层和周期性复苏网络。节肢族方案像一张暗红色神经网，连接未来生物群体，让个体意识可以随时并入集群。渊民族方案沉入海洋，蓝色声纹沿洋流扩散，像要让水本身成为记忆。震翼族方案最淡，只是一条条银色路线，连接陆地、海洋、天空和轨道。

没有哪一个方案完全丑陋。

也没有哪一个方案完全无害。

洛安看着那些叠在一起的未来，忽然觉得地球像被许多透明手掌按住。每只手都说自己只是想扶它一把，但如果所有手同时用力，种子会被捏碎。

“所以战争不是为了法厄同。”洛安说。

“法厄同只是门。”瑟恩说。

“门后是地球。”

“门后是未来。”渊民族老人纠正。

日辉族幸存者看向洛安：“没有引导，未来会浪费它自己。”

鳞裔族女人说：“没有保存，未来会毁在一次错误里。”

节肢族个体说：“没有并联，未来个体会重复孤独和冲突。”

瑟恩冷冷道：“没有路，未来只是你们几族轮流上锁的房间。”

所有人又一次沉默。

地球投影缓慢旋转。

这颗年轻行星没有为自己辩护。它只是转着，海洋反射着远古太阳的光，像还不知道自己会在亿万年后长出蓝鲸、人类、战争、梦，以及一个叫查理·黑尔的普通人。

巨鲸种子在洛安掌心轻轻震动。

我忽然觉得，它也在看地球。

不是像其他族群那样看一份方案。

而像看一个必须被保护的未完成答案。
